\section{The pull--map--push calculus}

An edge expression denotes a function from contexts to values.  The structural
fragment contains
\[
  \mathsf{sourceId},\quad \mathsf{destinationId},\quad \mathsf{edgeId}.
\]
Attribute expressions extend this fragment by precomposition with $s$, $d$, or
the edge identifier.  Thus a source attribute $a:V\to A$ denotes
$a\circ\mathsf{src}$, and similarly for destination and edge attributes.  This
is the \emph{pull} phase.

For an expression $e:C_G\to A$ and pure function $f:A\to B$, mapping and
emission have the denotation
\[
  \emit_G(F,e,f)=\mathsf{map}(f\circ e,\trav_G(F)).
\]
The implementation is not required to construct the intermediate context or
expression streams.

Let $M=(A,\oplus,0)$ be a monoid.  Source reduction produces one value for
every frontier occurrence:
\[
  \reduceSrc_G(F,e,f,M)[i]
  = \mathop{\bigoplus}_{c\in X_G(F[i])} f(e(c)).
\]
An isolated vertex therefore contributes $0$, and a repeated vertex produces
repeated output positions.

Destination reduction pushes values back to a dense vertex domain:
\[
  \reduceDst_G(F,e,f,M)[v]
  = \mathop{\bigoplus}_{\substack{c\in\trav_G(F)\\\mathsf{dst}(c)=v}}
      f(e(c)).
\]
Associativity and identity justify parallel tree reduction.  If the lowering
may reorder colliding messages, commutativity is additionally required for a
deterministic result.  Floating-point addition therefore needs either a weaker
equivalence relation or a lowering with an explicitly specified order; it is
not silently treated as an exact commutative monoid.
